\subsection{The remaining multiplicity-eight ring}
\label{subsec:frontier-1660}

The ring $R_{6,8,1660}$ is self-dual and contains the Ising based
subcategory with $g=b_1$ and $\sigma=b_4$. The unrestricted signed
reduction recalled above reduces every categorification to the
pseudo-unitary case. Its complete replay is
\filename{research/round5/m8_1660_conductor/replay_reduction.py}.
We use positive dimensions and put $D=\operatorname{FPdim}(\mathcal C)$.

\begin{proposition}[Actual-twist support and individual capacity]
\label{prop:1660-actual-support}
For a central simple with restriction $v=(v_0,\ldots,v_5)$, put
\[
 A(v)=v_0+v_1+v_2+v_3+2v_4-2v_5,\qquad
 B(v)=v_0+v_1-v_2-v_3+v_4+v_5.
\]
Then $A(v)\ne0$ implies $\theta_v^{24}=1$, and $B(v)\ne0$
implies $\theta_v^{12}=1$. Every actual central twist has order
dividing $48$.

Every high column containing $\sigma$ satisfies
\[
 v_0=v_1=0,\qquad v_4=1,\qquad v_2+v_3=4,\qquad v_5=3.
\]
There are exactly eight such central simples, each of dimension
$D\sqrt2/8$. Each of the four allowed actual Ising twists occurs on
exactly two of these simples.
\end{proposition}
\begin{proof}
The two saturated dimension-field characters have conductors $24$ and
$12$. Evaluating their Fourier coefficients in
\eqref{eq:saturated-column-support} gives nonzero constant multiples
of $A$ and $B$, respectively. The character representatives are
$1,5,7,11$; their sign vectors are $(1,1,-1,-1)$ and
$(1,-1,-1,1)$. The exact coefficient calculation is reconstructed by
\filename{software/dimension_block_support.py}.
Proposition~\ref{prop:saturated-conductor-support} gives the order bounds.

If a high column contains $\sigma$, its actual twist is primitive of
order $16$ and its $g$-coordinate is zero, by
Corollary~\ref{cor:fibonacci-ising-projection}. Therefore both Fourier
coefficients vanish. Solving these two integer linear equations gives
$v_2+v_3=4v_4$ and $v_5=3v_4$. Exact evaluation of the dimensions gives
\[
 \frac{v_4\operatorname{FPdim}(v)}{D\sqrt2/4}
       =\frac{v_4^2}{2}.
\]
The individual capacity bound \eqref{eq:individual-projection-capacity}
makes this ratio at most one. Hence $v_4=1$.
The induction Gram entry at $(\sigma,\sigma)$ is eight, and no low
column contains $\sigma$. Thus there are eight such simples, each
contributing half of one root mass, proving the stated allocation.
This uses individual actual-root masses before relative trace projection.

For the exponent bound, consider a high column on which both $A$ and
$B$ vanish. If $v_4>0$, the Ising support gives order $16$.
Otherwise $v_5=2(v_0+v_1)$ and $v_2+v_3=3(v_0+v_1)$.
Since $v_0=0$ and the column is nonzero, $v_1>0$, so its twist is $-1$.
All other high columns have twist order dividing $24$, and low twists
are one. Their least common multiple divides $48$.
\end{proof}

\begin{proposition}[Complete necessary restriction profiles]
\label{prop:1660-profiles}
For a categorification of $R_{6,8,1660}$, the pseudo-unitary Galois
conjugate has one of two low frames, each with two admissible character
labellings. After the actual-twist support and capacity bounds,
the complete induction Gram equations, joint-character and
normalized-dimension Galois balance, integer cubic Verlinde
multiplicities, and nonnegative low-product restriction remainders,
the numbers of high restriction profiles are as follows, per labelling:
\[
\begin{array}{c|r|r|c}
 \text{low frame}&\text{cubic constraints}&\text{words through degree six}
 &\text{range of possible center ranks}\\\hline
 0&371&371&25\text{--}76\\
 1&2011&1999&25\text{--}82
\end{array}
\]
For a fixed low frame the two labellings give the same surviving
restriction profiles. Every retained profile also passes entrywise
nonnegativity of the residual Gram matrix of all fifteen quadratic
words in the five nonunit low simples.
\end{proposition}
\begin{proof}
Enumerate the high columns using the induction Gram bounds,
formal-codegree divisibility, complete character integrality and
Galois closure, and the Ising projection constraints. Apply
Proposition~\ref{prop:1660-actual-support} and
Lemma~\ref{lem:simple-restriction-support}; the latter bounds the pure
$g$ multiplicity by one. There remain $20$ high columns in the first
frame and $52$ in the second.

Their nonnegative integral multiplicities satisfy a bounded affine
system. The Gram and Galois equations have ranks $11$ and $30$;
exact row reduction leaves $9$ and $22$ free coordinates. Diagonal
Gram entries bound every multiplicity. Cubic character formulas give
integrality congruences and nonnegativity inequalities. Restricting
each low product to the source based ring supplies coordinate bounds.

The bounded free coordinates are enumerated with exact interval and
divisibility pruning. The complete searches visit $1009$ and $63366$
nodes, respectively, per character labelling. Higher-word tests use
\eqref{eq:unified-word-pairing}; degrees four through six remove
twelve profiles in the second frame and none in the first.
The quadratic-word residual Gram test subtracts the known low
constituents as in Theorem~\ref{thm:joint-restriction-word-gram}.

All domains, rational row reductions, pruning decisions and retained
profiles are reconstructed by \filename{software/round5_1660.py}.
The records are \filename{research/round5/m8_1660/low0_profiles.json}
and \filename{low1_profiles.json} in the same directory.
The source-bound receipts \filename{low0_replay.json} and
\filename{low1_replay.json} record complete optimizer-free replays
under \texttt{python -O}. Consequently neither numerical solver
status nor Python assertions are used to prove completeness.
\end{proof}

The eight $\sigma$-supported simples have only two possible restriction
multisets in these lists. In frame zero all eight have column
$(0,0,2,2,1,3)$. In frame one, $1628$ profiles have this same multiset;
the other $371$ have six copies of this column and one each of
$(0,0,1,3,1,3)$ and $(0,0,3,1,1,3)$, per labelling.
The exact first-trace cancellation condition on opposite actual Ising
roots admits an allocation for both multisets, so it supplies no
additional exclusion by itself. The exact control
\filename{research/round5/m8_1660/verify_sigma_partitions.py}
reconstructs these multisets and their allocations, and checks the
displayed Fourier forms on all six basis columns.

\begin{proposition}[Complete actual-root profile reduction]
\label{prop:1660-actual-frontier}
The actual-root equations leave precisely the following necessary
profile alternatives:
\[
\begin{array}{c|c|r|r}
\text{low frame}&\text{character labelling}&\text{input profiles}
 &\text{remaining profiles}\\\hline
0&(0,5,4,1,3,2)&371&168\\
1&(0,5,4,1,2,3)&1999&798\\
1&(0,5,4,1,3,2)&1999&798
\end{array}
\]
All three branches have second indicators $(1,1,1,-1,1)$.
Every retained profile has eight copies of the balanced
$\sigma$-supported column $(0,0,2,2,1,3)$.
Their possible center ranks are $28,31,\ldots,76$ in the first frame
and $28,31,\ldots,82$ in the second.
\end{proposition}
\begin{proof}
The proved exponent bound permits the use of the $48$ actual roots as
the complete twist domain. Impose the induction Gram equations,
first and second indicator traces, cyclic-rotation multiplicities,
joint-character reciprocity, normalized-dimension Galois balance,
the four individual Ising-root masses, and the low-product
restriction constraints. There are sixteen low-frame, character,
and indicator branches. Exact Farkas vectors exclude thirteen of
them. In the first retained branch, three shared separating vectors
exclude $203$ of the $371$ restriction profiles.

For each second-frame branch the shared separating vectors leave
$966$ profiles. A six-row integer combination of the fixed-profile
system has every coefficient divisible by four, but its right-hand
side is not divisible by four for any asymmetric $(1,6,1)$
$\sigma$-pattern. This excludes all such patterns, including the
$168$ that remained after the separating-vector tests, leaving
$798$ profiles per labelling. The checker intersects these results
with the complete degree-six survivor indices; it does not truncate
the stored $2011$-element cubic list.

The compact record
\filename{research/round5/m8_actual48/certificate.json} stores the
shared integer vectors, congruences, and complete remaining index
lists. Its replay program \filename{certificate.py} reconstructs
all sixteen systems and checks every applicable profile exactly.
The source-bound \filename{replay_O.json} records a complete run
with numerical optimizers and Python assertions disabled.
\end{proof}

\begin{proposition}[A forced Ising-center subrepresentation]
\label{prop:1660-ising-subrepresentation}
The canonical modular representation of a hypothetical center
contains the full nine-dimensional representation of the Ising
center as a subrepresentation. Its central projection matrix $B$
satisfies $BB^\dagger=2I_9$.
\end{proposition}
\begin{proof}
For every occupied column in the complete lists of
Proposition~\ref{prop:1660-profiles}, its Ising projection is one of
\[
 (0,0,0),\quad(1,0,0),\quad(1,1,0),\quad(0,1,0),\quad(0,0,1).
\]
At the allowed actual twist, each nonzero projection is a single
simple of the Ising center, with multiplicity one. This finite
support assertion is reconstructed for all $4740$ labelled profiles
by \filename{research/round5/m8_modrep/ising_complement.py}.
Corollary~\ref{cor:disjoint-central-projection} therefore applies:
the rows of $B$ have disjoint supports, their individual positive
dimension masses make them nonzero, and $BB^\dagger$ is a positive
scalar matrix. The four $\sigma$ rows each contain exactly two
objects, so that scalar is two.
\end{proof}

The finite-quotient criterion of
Subsection~\ref{subsec:modular-word-closure} is implemented for
$\mathrm{SL}_2(\mathbb Z/48\mathbb Z)$ using its two Chinese-remainder
factors. The exact class algebras give $76$ irreducible characters
for $\mathrm{SL}_2(\mathbb Z/16\mathbb Z)$ and seven for
$\mathrm{SL}_2(\mathbb Z/3\mathbb Z)$. Their $532$ tensor products
give $504$ distinct signed $T$-spectra. The independent replay checks
all class-algebra identities, character orthogonality, degrees,
central signs and exact Fourier multiplicities. Two Ising-center
spectra serve as positive controls. These records are in
\filename{research/round5/m8_modrep/}.

One may apply this spectrum cone to the complement of the forced
Ising-center representation. Its known $T$-spectrum has multiplicity
three at $1$, two at $-1$, and one at each of the four Ising roots.
The signed-trace lower bounds must also subtract these multiplicities:
they become $3$ at $1$, $-2$ at $-1$, and $-1$ at the four Ising roots.
The complement itself need not have permutation duality.

The full-spectrum and forced-complement searches produced no further
exact exclusion. The finite-group character enumeration is complete;
the coupled nonnegative integral completion problem for the remaining
profiles is not solved. The ring remains unresolved. The remaining
lists are necessary alternatives, not induction matrices or modular
data asserted to come from fusion categories.
